\section{\methodname: A unified taxonomy of visual capabilities}

Existing benchmarks typically organize visual tasks by task formulation or output modality, leaving visual generation and understanding tasks separated.
This makes it difficult to compare and study tasks that rely on similar visual capabilities but produce different outputs.
We introduce \methodname{}, which organizes I2I tasks and understanding tasks in a shared hierarchy.
For a visual understanding sample, we consider the primary visual capability required to solve it; for an I2I task, we consider the visual capability directly supervised by its training objective.
At the top level, we adopt the three Rs of computer vision~\citep{malik2016three}: \textit{Recognition}, \textit{Reconstruction}, and \textit{Reorganization}.
I2I training objectives and understanding capabilities remain separate leaves but are placed in the same hierarchy according to these shared visual capabilities.

\begin{figure}[t]
\centering
\includegraphics[width=\linewidth]{iclr2026/figures/OmniTaskonomy-standalone.pdf}
\caption{\textbf{\methodname.}
A unified taxonomy of visual tasks organized into three broad families: \textit{Recognition}, \textit{Reconstruction}, and \textit{Reorganization}.
I2I training objectives and understanding capabilities remain separate leaves, but share the same visual hierarchy.}
\label{fig:taskonomy}
\end{figure}

\paragraph{Organizing visual capabilities with the 3Rs.}
\textbf{Recognition} covers 
% \ranjay{concerns is weird english}
semantic content, including object and part identity, appearance, state, activity, text, and symbols.
\textbf{Reconstruction} covers geometric and photometric properties of the scene, including spatial relations, depth, metric distance, orientation, 3D shape, lighting, and occlusion.
\textbf{Reorganization} covers how visual elements are grouped, localized, and related, including segmentation, correspondence, ordering, counting, keypoints, and connectivity.
The 3R families define the top level of \methodname{}; fine-grained capabilities are derived from existing visual tasks rather than specified by the three families.
Appendix~\ref{app:unitaskonomy-definitions} lists all understanding capabilities and their definitions.

\paragraph{Deriving fine-grained visual capabilities.}
We derive the fine-grained visual understanding capabilities bottom-up from seven vision-language benchmarks: BLINK~\citep{fu2024blink}, MMStar~\citep{chen2024wemmstar}, MMT-Bench~\citep{ying2024mmt}, CV-Bench~\citep{tong2024cambrian}, RealWorldQA~\citep{xai2024realworldqa}, MMVP~\citep{tong2024eyes}, and VStarBench~\citep{wu2024vstarbench}.
We first sample benchmark questions and use \texttt{gemini-3-flash-preview}~\citep{gemini3flash} to describe the visual attributes and relations involved in solving each question, such as \emph{relative depth}, \emph{size comparison}, and \emph{occlusion}.
Then we review these descriptions and consolidate them into a fixed attribute schema, a list of key--value pairs describing image attributes, which is then used to annotate the full benchmark collection so that tasks from different benchmarks are described consistently.
We next divide the annotated samples into ten benchmark-stratified batches.
For each batch, a VLM organizes the samples into a local task tree using their annotated attributes and assigns each sample to the leaf that best describes its primary visual requirement.
We then merge the local trees, combining synonymous leaves across batches to form an initial set of fine-grained capabilities.
We place the resulting capabilities under \textit{Recognition}, \textit{Reconstruction}, or \textit{Reorganization}.
We manually review the resulting hierarchy, splitting leaves that mix distinct capabilities, merging redundant leaves, and refining definitions until the leaf categories have clear boundaries.
Figure~\ref{fig:taskonomy} shows the resulting taxonomy.
Appendix~\ref{app:construction-schema} describes the attribute schema, and Appendix~\ref{app:construction-trees} provides further details on tree construction and refinement.

\paragraph{Integrating I2I tasks.}
We next place I2I training objectives in the same 3R hierarchy.
We collect a broad set of dense prediction, reconstruction, and image-editing tasks from existing vision and multimodal learning settings, including tasks studied in Taskonomy~\citep{zamir2018taskonomy} as well as objectives such as inpainting and object editing.
Each I2I task is placed according to the main visual capability directly supervised by its objective.
For example, Z-depth and Euclidean depth prediction are placed under \textit{Reconstruction}, alongside the understanding capability \textit{depth understanding}.
2D keypoint prediction is placed under \textit{Reorganization}.
Object editing and attribute editing are placed under \textit{Recognition}, as they directly supervise object identity and visual attributes such as color or material.
The final taxonomy contains 19 I2I tasks and 25 understanding capabilities.
Appendix~\ref{app:unitaskonomy-i2i} lists all I2I tasks and their assignments, and Appendix~\ref{app:i2i-training-data} describes their training data.

\paragraph{Annotating benchmark samples.}
Once the taxonomy is fixed, we reassign each eligible benchmark sample to one of the finalized understanding leaves.
Three VLM judges independently annotate each sample using the image, question, answer choices, reference answer, and the same capability definitions.
We retain an example when at least two of the three judges assign it to the same capability and exclude examples without a majority, yielding 9,444 samples.
Appendix~\ref{app:unitaskonomy-annotation} reports unanimous agreement, majority agreement, and exclusion statistics.
Appendix~\ref{app:prompt-annotation} provides the annotation prompts and model inputs.
We further validate the resulting assignments through human reviews.
Four human annotators review a stratified subset of 250 samples covering all 25 capabilities.
Each annotator reviews 50 examples shared across all four annotators and 50 additional non-overlapping examples, yielding 400 human reviews.
For each example, the annotator can accept the proposed capability, reassign it to another capability, or mark the assignment as unsure.
Among definite judgments, 97.4\% of human labels match the VLM-assigned capability.
Human-VLM agreement remains high after correcting for chance, with Cohen's $\kappa$ ranging from 0.956 to 0.989 across annotators (mean $0.972$).
On the 50 examples reviewed by all four annotators, the resulting human capability labels also show high inter-annotator agreement (Krippendorff's $\alpha=0.943$).
These results provide human validation of the automated sample assignments.
Appendix~\ref{app:human-validation} gives the full review protocol and agreement analysis.
